\documentclass[../../../include/open-logic-section]{subfiles}

\begin{document}

% Part: incompleteness
% Chapter: incompleteness-provability
% Section: fixed-point-lemma

\olfileid{inc}{inp}{fix}

\olsection{不动点引理}

\begin{explain}
不动点引理表明，只要 $\Th{T}$ 扩张了~$\Th{Q}$，对于任意公式~$!B(x)$，都存在一个语句~$!A$ 使得 $\Th{T} \Proves !A \liff !B(\gn{!A})$。
在说谎者语句的情形中，我们想要 $!A$（在~$\Th{T}$ 中可证地）等价于“$\gn{!A}$ 是假的”，
也就是说，等价于“$\Gn{!A}$ 是一个假语句的Gödel数”这一陈述。
为了理解证明的思路，不妨将其与 Quine对~$!A$ 的非形式解读作一比较：
$!A$ 即“‘当冠以其自身的引述时产生假句’当冠以其自身的引述时产生假句”。
取出一个表达式，然后将该表达式的自身引述冠于其前从而构成一个语句，
这种操作可以称为对该表达式的\emph{对角化}，其结果称为它的对角化。
所以，“当冠以其自身的引述时产生假句”的对角化就是
“‘当冠以其自身的引述时产生假句’当冠以其自身的引述时产生假句”。
请注意，Quine 的说谎者语句并不是“产生假句”的对角化，
而是“当冠以其自身的引述时产生假句”的对角化。
因此，被对角化以产生说谎者语句的那个性质本身就已经涉及对角化！

在算术语言中，我们通过计算含一个自由变元的公式的Gödel数，
然后将代表该Gödel数的标准数字代入自由变元，来形成该公式的引述。
$!E(x)$ 的对角化是 $!E(\num{n})$，其中 $n = \Gn{!E(x)}$。
（从现在起，我们把 $\num{\Gn{!E(x)}}$ 简写为 $\gn{!E(x)}$。）
所以，如果 $!B(x)$ 是“是假句”，那么
“当冠以其自身的引述时产生假句”就可以表达为
“当应用于其对角化的Gödel数时产生假句”。
如果我们有一个符号~$\Obj{diag}$ 代表函数 $\fn{diag}(n)$
——该函数计算Gödel数为~$n$ 的公式的对角化的Gödel数——那么我们就可以将 $!E(x)$ 写成
$!B(\Obj{diag}(x))$。
而 Quine 版本的说谎者语句就将是它的对角化，即 $!E(\gn{!E(x)})$ 或
$!B(\Obj{diag}(\gn{!B(\Obj{diag}(x))}))$。
当然，$!B(x)$ 可以是任何其他性质，同样的构造仍能奏效。
对于不完全性定理，我们将取 $!B(x)$ 为“$x$ 在~$\Th{T}$ 中不可推导”。
那么 $!E(x)$ 就是“当应用于其对角化的Gödel数时，产生一个在~$\Th{T}$ 中不可推导的语句”。

为了在~$\Th{T}$ 中将此形式化，我们必须找到将 $\fn{diag}$ 形式化的方法。
函数 $\fn{diag}(n)$ 是可计算的，事实上，它是原始递归的：
如果 $n$ 是公式~$!E(x)$ 的Gödel数，
$\fn{diag}(n)$ 就返回 $!E(\gn{!E(x)})$ 的Gödel数。
（回想一下，$\gn{!E(x)}$ 是~$!E(x)$ 的Gödel数的标准数字，即 $\num{\Gn{!E(x)}}$。）
如果 $\Obj{diag}$ 是 $\Th{T}$ 中代表函数 $\fn{diag}$ 的函数符号，
我们就可以取 $!A$ 为公式 $!B(\Obj{diag}(\gn{!B(\Obj{diag}(x))}))$。注意到
\begin{align*}
\fn{diag}(\Gn{!B(\Obj{diag}(x))}) & =
\Gn{!B(\Obj{diag}(\gn{!B(\Obj{diag}(x))}))} \\
& = \Gn{!A}.
\end{align*}
假设 $\Th{T}$ 可以推导出
\[
\Obj{diag}(\gn{!B(\Obj{diag}(x))}) = \gn{!A},
\]
它就能推导出 $!B(\Obj{diag}(\gn{!B(\Obj{diag}(x))}))
\liff !B(\gn{!A})$。而左边根据定义正是~$!A$。

当然，$\Obj{diag}$ 一般不会是 $\Th{T}$ 的函数符号，更不是~$\Th{Q}$ 的。
但是，既然 $\fn{diag}$ 是可计算的，它在~$\Th{Q}$ 中就是\emph{可表示的}，即可由某个公式
$!D_{\fn{diag}}(x,y)$ 表示。
因此，我们可以把 $!B(\Obj{diag}(x))$ 写成 $\lexists[y][(!D_{\fn{diag}}(x,y) \land !B(y))]$。
其余部分，上面概述的证明依然成立，而且事实上，它在~$\Th{Q}$ 中就已经成立。
\end{explain}

\begin{lem}
\ollabel{lem:fixed-point} 设 $!B(x)$ 是任意一个含一个自由变元~$x$ 的公式。
则存在语句~$!A$，使得 $\Th{Q} \Proves
!A \liff !B(\gn{!A})$。
\end{lem}


\begin{proof}
给定 $!B(x)$，令 $!E(x)$ 为公式
$\lexists[y][(!D_{\fn{diag}}(x,y) \land !B(y))]$，并令 $!A$~为其对角化，即公式 $!E(\gn{!E(x)})$。

由于 $!D_{\fn{diag}}$ 表示 $\fn{diag}$，且
$\fn{diag}(\Gn{!E(x)}) = \Gn{!A}$，$\Th{Q}$ 可以推导出
\begin{align}
  & !D_{\fn{diag}}(\gn{!E(x)}, \gn{!A}) \ollabel{repdiag1} \\
  & \lforall[y][(!D_{\fn{diag}}(\gn{!E(x)},y) \lif
  \eq[y][\gn{!A}])]. \ollabel{repdiag2}
\end{align}
现在我们证明 $\Th{Q} \Proves !A \liff !B(\gn{!A})$。我们非形式地进行论证，仅使用逻辑和在~$\Th{Q}$ 中可推导的事实。

首先，假设~$!A$，即 $!E(\gn{!E(x)})$。回到 $!E(x)$ 的定义，
我们看到 $!E(\gn{!E(x)})$ 正好就是
\[
\lexists[y][(!D_{\fn{diag}}(\gn{!E(x)},y) \land !B(y))].
\]
考虑这样一个~$y$。由于 $!D_{\fn{diag}}(\gn{!E(x)},y)$，根据
\olref{repdiag2}，有 $y = \gn{!A}$。所以，从 $!B(y)$ 可得~$!B(\gn{!A})$。

现在假设 $!B(\gn{!A})$。根据 \olref{repdiag1}，我们有
\begin{align*}
& !D_{\fn{diag}}(\gn{!E(x)}, \gn{!A}) \land !B(\gn{!A}).
\intertext{由此可得}
& \lexists[y][(!D_{\fn{diag}}(\gn{!E(x)},y) \land !B(y))].
\end{align*}
但这正是 $!E(\gn{!E(x)})$，即~$!A$。
\end{proof}

\begin{digress}
你应该将此与可计算性理论中不动点引理的证明作比较。
区别在于，这里我们想用自身来定义一个\emph{语句}，而那里我们想用自身来定义一个\emph{函数}；
撇开这一差异，两者的思想确实是一样的。
\end{digress}

\begin{prob}
  如果 $\Th{Q}
  \Proves !B \liff !A(\gn{!B})$ 对所有的语句~$!B$ 都成立，则称公式~$!A(x)$ 是一个\emph{真值定义}。试借助不动点引理证明，不存在真值定义公式。
\end{prob}

\end{document}